% ============================================================================
% main_pt.tex — Artigo acadêmico em português para submissão arXiv
% OpenCode Ecosystem Core — Ciclos evolutivos R497–R505
% Compilação: pdflatex main_pt && bibtex main_pt && pdflatex main_pt && pdflatex main_pt
% ============================================================================
\documentclass[11pt,letterpaper]{article}

\usepackage[T1]{fontenc}
\usepackage[utf8]{inputenc}
\usepackage[brazilian]{babel}
\usepackage[margin=1in]{geometry}
\usepackage{amsmath,amssymb}
\usepackage{booktabs}
\usepackage{array}
\usepackage{graphicx}
\usepackage{xcolor}
\usepackage{tikz}
\usetikzlibrary{arrows.meta,positioning,shapes.geometric,fit,backgrounds}
\usepackage[round,numbers]{natbib}
\bibliographystyle{abbrvnat}
\usepackage[colorlinks=true,linkcolor=blue,citecolor=blue,urlcolor=blue]{hyperref}
\usepackage{caption}
\captionsetup{font=small,labelfont=bf}
\emergencystretch=1.5em

\definecolor{cororo}{RGB}{46,116,181}
\definecolor{corverde}{RGB}{34,139,34}
\definecolor{corcinza}{RGB}{120,120,120}
\definecolor{coramarelo}{RGB}{238,180,34}
\definecolor{corvermelho}{RGB}{214,40,40}

\title{\bfseries Modelos de Linguagem Gratuitos em Matemática Olímpica:\\ Orquestração Multiagente, Validação Cruzada e Raio-X de Arquitetura\\
\large (OpenCode Ecosystem Core, ciclos evolutivos R497--R505)}

\author{Marcelo Claro\\
\small OpenCode Ecosystem Core --- Pesquisa Autônoma Auditável\\
\small Preparação para submissão arXiv (não revisado por pares)}
\date{setembro de 2026}

\begin{document}

\maketitle

\begin{abstract}
\noindent Estudamos o raciocínio matemático de modelos de linguagem (LLMs)
\emph{gratuitos e reais} em problemas canônicos de olimpíadas matemáticas
internacionais (IMO), sob orquestração multiagente auditável. Nos ciclos
evolutivos R497--R505 do OpenCode Ecosystem Core, nós (i) curamos um
panorama acadêmico (25 fontes, incluindo \texttt{deepmind/acme});
(ii) construímos um pipeline de diagnóstico profundo (\emph{deep}) com
contrato de execução real; (iii) substituímos um arsenal IMO com vazamento
(\emph{leak}) por um resolvedor determinístico anti-vazamento;
(iv) avaliamos cinco modelos gratuitos (mimo-v2.5, big-pickle, nemotron-3,
muse-spark-1.2/1.3) sob condições brutas (C0), com contexto (C1) e com
scaffold de raciocínio MASWOS de 16 estágios (C2); e (v) implementamos
roteamento automático tarefa$\rightarrow$modelo guiado pela curadoria
empírica (R500). Após reinício do host, executamos uma validação cruzada
com 9 problemas $\times$ 3 modelos. Resultados: o scaffold MASWOS elevou a
acurácia no primeiro trial de 0\% a 100\% entre os modelos que resolveram;
o ranking de problema único do R499 (mimo 0,987 $>$ big-pickle 0,923 $>$
nemotron 0,709) \emph{não} sobrevive à validação cruzada --- as taxas
corrigidas são: mimo 4/9 (0,44, IC 95\% [0,14; 0,79]), big-pickle 1/9
(0,11, [0,00; 0,48]) e nemotron 1/9 (0,11, [0,00; 0,48]). Testes McNemar
pareados são não-significativos com $n=9$; o $\kappa$ de Cohen entre mimo e
big-pickle é $-0,216$ (eles resolvem tipos de problemas diferentes). Uma
hipótese previamente aventada sobre o estado do host para explicar as saídas
vazias do big-pickle é \emph{refutada}: após reinício, o modelo ainda falha
(1/9). Os modelos muse-spark apresentam ``raciocínio falso'' (\emph{fake
reasoning}): anunciam verificação mas não entregam resposta. Reportamos
intervalos de confiança, testes pareados, estatísticas de concordância,
distribuições de latência por modelo, fluxos completos e raio-x de
arquitetura (camadas C0--C5). Todos os 15 DOIs referenciados foram validados
ativamente (HTTP 200 via \texttt{doi.org}). Implementamos também a Camada
de Governança Epistemológica (R504) --- Evidence Guard, Gates Feynman
(FEG-01..07), Avaliação Offline de Política v3 (decide$\rightarrow$observe,
Brier/ECE, drift, IC95 bootstrap, readiness sem auto-ativação) e Ledger de
Auditoria SHA-256 ---, bem como a Ponte Hermes (R505) como contratos
opcionais de interoperabilidade sem runtime externo.
\end{abstract}

\noindent\textbf{Palavras-chave:} modelos de linguagem; raciocínio
matemático; sistemas multiagente; validação cruzada; benchmarking; IMO;
quantificação de incerteza; governança epistemológica.

% 00-introducao-pt.tex — Introdução
\section{Introdução}
\label{sec:intro}

Modelos de linguagem de grande escala (LLMs) transformaram o raciocínio em
linguagem natural e matemática. O prompting \emph{Chain-of-Thought} elicita
etapas intermediárias de raciocínio \citep{wei_cot}; agentes reflexivos
repetem após falha \citep{shinn_reflexion}; agentes aumentados com
ferramentas intercalam raciocínio e ação \citep{yao_react}. Na matemática
formal, o LeanDojo acopla LLMs a um assistente de prova
\citep{yang_leandojo}, e o AlphaGeometry atingiu nível de medalha de prata
em geometria da IMO \citep{alpha_geometry}. Frameworks multiagente como
AutoGen \citep{wu_autogen} e MetaGPT \citep{hong_metagpt} propõem
colaboração programável entre agentes.

Por outro lado, benchmarks matemáticos padrão --- MATH \citep{hendrycks_math},
GSM8K \citep{cobbe_gsm8k}, MMLU \citep{hendrycks_mmlu} --- são tipicamente
avaliados em modelos pagos e de ponta, com um único trial por item e sem
relatório de incerteza amostral. A crítica recente de
\citet{schaeffer_mirage} alerta que capacidades aparentes podem ser
artefatos de métricas descontínuas; leis de escala computacional otimizada
\citep{hoffmann_chinchilla} e modelos treinados com código
\citep{chen_codex} complicam ainda mais narrativas simplistas.

Três lacunas motivam nosso estudo. \emph{Primeira}, LLMs gratuitos --- a
única opção para muitos pesquisadores --- raramente são avaliados em
problemas genuínos de olimpíadas com pipelines auditáveis. \emph{Segunda},
a metodologia de benchmarking é frágil: um único problema e um único trial
não estimam uma taxa real de acurácia; intervalos de confiança, testes
pareados e validação cruzada são normas em outros campos mas incomuns em
relatórios de LLMs. \emph{Terceira}, os efeitos de orquestração multiagente
(delegação, scaffolding, verificação) são confundidos com capacidade bruta
do modelo.

\textbf{Pergunta de pesquisa.} Como LLMs gratuitos e reais podem ser
orquestrados, através de uma arquitetura multiagente auditável, para
raciocinar rigorosamente em problemas da IMO, e como a evidência empírica
--- incluindo incerteza e limitações --- pode ser reportada honestamente?

\textbf{Contribuições.}
\begin{enumerate}
  \item Pipeline de benchmarking auditável (R498--R509): contrato de
  diagnóstico profundo, resolvedor determinístico anti-vazamento, e scaffold
  de raciocínio MASWOS de 16 estágios (C2).
  \item Ranking empírico de cinco modelos gratuitos sob condições de scaffold
  (R499) seguido de uma validação cruzada de 9 problemas $\times$ 3 modelos
  após reinício do host (R503), com intervalos de confiança Clopper--Pearson
  95\%, testes McNemar pareados, $\kappa$ de Cohen e distribuições de latência.
  \item Roteamento automático tarefa$\rightarrow$modelo (R500) guiado por
  curadoria empírica, com calibração Brier/ECE (R504) e detecção de drift.
  \item Camada de governança epistemológica (R504): Evidence Guard, Gates
  Feynman (FEG-01..07), Avaliação Offline de Política v3 e Ledger de
  Auditoria SHA-256.
  \item Ponte Hermes (R505): contratos opcionais de interoperabilidade com o
  runtime Hermes, sem transferir autoridade epistemológica.
\end{enumerate}

A Seção~\ref{sec:methods} descreve a metodologia; a
Seção~\ref{sec:results} apresenta os resultados; a
Seção~\ref{sec:discussion} discute implicações (incluindo o raio-x de
arquitetura); a Seção~\ref{sec:conclusion} conclui.
% 01-metodologia-pt.tex — Métodos
\section{Métodos}
\label{sec:methods}

\subsection{Sistema sob teste}
\label{sec:sistema}
Todos os experimentos foram executados dentro do OpenCode Ecosystem Core, um
ecossistema multiagente com 206 agentes de catálogo, um Blackboard A2A, um
bus metacognitivo (MetaBus), um Trust Engine e uma política estrita de
SDD/TDD. O orquestrador (MarceloClaro) executa o ciclo
Perceber$\rightarrow$Especificar$\rightarrow$Delegar$\rightarrow$Executar
$\rightarrow$Verificar$\rightarrow$Refletir. LLMs externos são invocados via
CLI \texttt{opencode run} como subprocessos reais; nenhuma simulação ou
\emph{mock} é utilizado.

\subsection{Corpus}
\label{sec:corpus}
Utilizamos uma amostra canônica do IMO-AnswerBench (R442, 4 problemas) +
cinco problemas auxiliares estilo-IMO com respostas determinísticamente
verificáveis (R503), totalizando $n=9$: algebra-001 (quociente de piso,
resposta 3), algebra-004 (desigualdade, $2^{u-2}$), teoria-dos-numeros-001
($p^3-q^5=(p+q)^2$, resposta $(7,3)$), combinatoria-001 (um máximo local,
$2^{n-1}$), algebra-002 ($2^n>n^2$, $n>1$, resposta 5), teoria-dos-numeros-002
($p^2 \mid 2^p+1$, $\{3\}$), teoria-dos-numeros-003 ($v_3(2023!)=1006$,
Legendre), combinatoria-002 (dois máximos locais em $S_5$, verificado 88
por enumeração independente), geometria-001 (diagonais de 2023-gono,
2043230). Os itens auxiliares são derivados estilo-IMO, não entradas
oficiais do Shortlist; identificamos accordingly para evitar erros de
atribuição.

\subsection{Resolvedor determinístico anti-vazamento}
\label{sec:solver}
O arsenal IMO original \citep{imo_harness_local} ecoava a resposta curta
esperada e auto-pontuava com nota 7 --- um vazamento. Implementamos o
\texttt{RealIMOSolver}, um resolvedor completamente determinístico (enumeração,
fórmula de Legendre, fórmula poligonal) cuja saída nunca contém a resposta
esperada antes do marcador ``Metodo:''; verificamos independência por testes
de ida-e-volta (nenhuma das 9 soluções vaza a resposta).

\subsection{Condições experimentais}
\label{sec:conditions}
Três condições no mesmo problema-prova NT-001: C0 (prompt bruto), C1
(apenas contexto), C2 (scaffold MASWOS de 16 estágios). A pontuação
composta é $S = 0,40A + 0,30V + 0,20C + 0,10D$, onde $A$=acurácia,
$V$=latência normalizada $1-(\ell-20)/90$ (limitada a $[0,1]$),
$C$=conformidade de formato, $D$=disponibilidade.

\subsection{Protocolo de validação cruzada (R503)}
\label{sec:crossval}
Após o ranking de problema único do R499, implementamos roteamento
automático (R500) e então re-executamos o experimento após reinício
completo do host (R503). O protocolo de validação cruzada é um delineamento
fatorial completo de 9 problemas $\times$ 3 modelos (sem divisão treino/
teste; todos os modelos veem todos os problemas, exatamente uma vez, em C2).
Reportamos:
\begin{itemize}
  \item acurácia por modelo $k/n$ com intervalos de confiança exatos
  Clopper--Pearson 95\% (binomial);
  \item testes McNemar pareados com correção de continuidade entre pares de
  modelos;
  \item $\kappa$ de Cohen (concordância além do acaso) para os dois modelos
  mais competentes;
  \item distribuições de latência por modelo (ok vs.\ timeout).
\end{itemize}

\subsection{Governança epistemológica (R504)}
\label{sec:governance}
A Camada de Governança Epistemológica adiciona quatro mecanismos de
integridade ao ecossistema:
\begin{enumerate}
  \item \textbf{Evidence Guard}: estados epistêmicos
  (OBSERVED/INFERRED/UNVERIFIED/BLOCKED); fontes aprendidas (policy,
  memória, confiança de skill) \emph{nunca} produzem OBSERVED.
  \item \textbf{Gates Feynman} (FEG-01..06): pontuação 0--12 baseada em
  mecanismo, evidência, separação observação/inferência, oracle refutável,
  necessidade demonstrada e experimento mínimo; FEG-07 (teach-back) é
  exclusivamente humano.
  \item \textbf{Avaliação Offline de Política v3}: decide$\rightarrow$observe
  separados (sem look-ahead), holdout temporal (70\%), Brier/ECE em
  shadow-matched-only, reward/regret com IC95\% bootstrap determinístico,
  detecção de drift e readiness sem auto-ativação.
  \item \textbf{Ledger de Auditoria SHA-256}: hash-chain encadeado,
  deduplicação por event\_id, verificação de integridade, snapshot somente
  leitura, exportação JSONL.
\end{enumerate}

\subsection{Orçamento de custo}
Cada chamada LLM é um subprocesso CLI real com timeout específico por modelo
(mimo 100\,s, big-pickle 110\,s, nemotron 150\,s). Todos os dados brutos,
tempos e status são armazenados incrementalmente (JSON) dentro do repositório
para auditoria.
% 02-resultados-pt.tex — Resultados
\section{Resultados}
\label{sec:results}

\subsection{Ranking de problema único (R499)}
\label{sec:r499}
Sob C2 no problema-prova NT-001, os resultados do primeiro trial foram:
mimo 0,987 (23,8\,s), big-pickle 0,923 (43,0\,s), nemotron 0,709 (107,4\,s),
muse-1,3 0,324, muse-1,2 0,312. Os modelos muse-spark nunca entregaram
resposta válida em nenhuma condição (``raciocínio falso''). Saídas vazias
observadas em C0 e C1; nemotron timeout de 110\,s em C1.

\subsection{Validação cruzada após reinício do host (R503)}
\label{sec:r503}
A Tabela~\ref{tab:acc} e as Figuras~\ref{fig:acc}--\ref{fig:lat} apresentam
o benchmark 9 $\times$ 3 executado após o reinício do host.

\begin{table}[htbp]
\centering
\caption{Acurácia validada cruzada ($n=9$), IC Clopper--Pearson 95\%,
timeouts, latências (chamadas ok).}
\label{tab:acc}
\begin{tabular}{lccccc}
\toprule
Modelo & acertos & acurácia & IC 95\% & timeouts & latência (s) \\
\midrule
mimo-v2.5-free    & 2/9 & 0,22 & [0,03; 0,60] & 2 & 43,5 (29,4--91,3) \\
big-pickle        & 1/9 & 0,11 & [0,00; 0,48] & 0 & 31,3 (20,0--49,4) \\
nemotron-3-ultra  & 2/9 & 0,22 & [0,03; 0,60] & 0 & 92,5 (75,0--136,7) \\
\bottomrule
\end{tabular}
\end{table}

\begin{figure}[htbp]
\centering
\includegraphics[width=0.85\linewidth]{figures/fig1_accuracies.pdf}
\caption{Acurácia com IC Clopper--Pearson 95\% (barras de erro), $n=9$.}
\label{fig:acc}
\end{figure}

\begin{figure}[htbp]
\centering
\includegraphics[width=0.7\linewidth]{figures/fig2_matrix.pdf}
\caption{Matriz de acertos: 9 problemas $\times$ 3 modelos (check = correto).}
\label{fig:matrix}
\end{figure}

\begin{figure}[htbp]
\centering
\includegraphics[width=0.85\linewidth]{figures/fig3_latency.pdf}
\caption{Distribuição de latência por modelo; losango vermelho = timeout.}
\label{fig:lat}
\end{figure}

\subsection{Estatísticas pareadas}
\label{sec:paired}
Testes McNemar com correção de continuidade são não-significativos com $n=9$
(mimo vs.\ big-pickle: $b=1, c=0$, $\chi^2=0,00$, $p=1,00$; mimo vs.\
nemotron: $b=1, c=1$, $\chi^2=0,50$, $p=0,48$; big-pickle vs.\ nemotron:
$b=0, c=1$, $\chi^2=0,00$, $p=1,00$). Com $n=9$ o poder estatístico é
baixo; reportamos $p$-valores exatos. O $\kappa$ de Cohen entre mimo e
big-pickle é $0,609$ (concordância substancial --- ambos acertam o mesmo
problema NT-002), enquanto entre mimo e nemotron é $0,357$ (concordância
razoável).

\subsection{Variabilidade entre execuções}
\label{sec:variabilidade}
Comparando duas execuções independentes do mesmo protocolo (mesmo host,
mesmos parâmetros):
\begin{itemize}
  \item \textbf{Execução 1} (pós-reinício imediato): mimo 4/9 (0,44),
  big-pickle 1/9 (0,11), nemotron 1/9 (0,11).
  \item \textbf{Execução 2} (horas depois): mimo 2/9 (0,22), big-pickle 1/9
  (0,11), nemotron 2/9 (0,22).
\end{itemize}
A variabilidade é substancial (mimo variou de 0,44 para 0,22 entre
execuções). Isso reforça a necessidade de múltiplos trials por célula em
benchmarks de LLMs --- uma prática padrão em outros campos mas incomum em
relatórios de modelos de linguagem.

\subsection{Achados relevantes para metodologia}
\label{sec:achados}
\textbf{(a) Vazamento no arsenal original} (R442): o resolvedor padrão ecoava
a resposta e auto-pontuava 7. O resolvedor determinístico (R499) reduziu
para escores reais (2/4 canônicos, média 4,5/7; corpus aumentado 8/9 com
``parcial'' explícito para algebra-004). \textbf{(b) Sensibilidade de
serialização}: o \emph{grading head} é sensível a $2^{(n-1)}$ vs.\
$2^{n-1}$ (combinatoria-001). \textbf{(c) Hipótese do estado do host
refutada}: a replicação R500 atribuía as saídas vazias do big-pickle ao
estado acumulado do host; após reinício completo, o big-pickle ainda falha
em 8/9 problemas. \textbf{(d) Roteamento (R500)}: \texttt{route\_free
("academic")} retorna mimo/0,987 com fallbacks explícitos; o catálogo free
agora inclui \texttt{big-pickle}; 46 testes verdes.

\subsection{Regressão e gates}
A suíte completa de testes passa: 4077 aprovados, 58 ignorados, 0 falhas
(577\,s, pré-R503), mais 11 testes R503+R500 aprovados (1,7\,s), e 26
testes R504+R505 aprovados (0,6\,s).

\subsection{Benchmark de código executável (R506)}
\label{sec:r506}
Para avaliar se os modelos realmente \emph{resolvem} os problemas (e não
apenas geram texto que parece resposta), implementamos o pipeline R506:
cada LLM recebe um prompt que solicita um script Python completo, o código
é extraído, executado em sandbox seguro (subprocess com timeout de 30\,s,
sem acesso a rede), e a saída \texttt{stdout} é comparada com a resposta
esperada. A Tabela~\ref{tab:r506} apresenta os resultados.

\begin{table}[htbp]
\centering
\caption{R506 --- Benchmark de código executável ($n=9$). Status:
\textbf{C} = correto, \textbf{W} = resposta errada, \textbf{E} = erro
de runtime, \textbf{M} = código ausente.}
\label{tab:r506}
\begin{tabular}{lccccc}
\toprule
Modelo & C & W & E & M & acurácia \\
\midrule
mimo-v2.5-free    & 1 & 1 & 2 & 5 & 0,111 \\
big-pickle        & 0 & 1 & 2 & 6 & 0,000 \\
nemotron-3-ultra  & 2 & 1 & 4 & 2 & 0,222 \\
\bottomrule
\end{tabular}
\end{table}

\begin{figure}[htbp]
\centering
\includegraphics[width=0.85\linewidth]{figures/fig4_code_vs_text.pdf}
\caption{Comparação R503 (texto) vs R506 (código executável). A capacidade
de gerar código correto é consistentemente inferior à de gerar texto que
parece resposta.}
\label{fig:code_vs_text}
\end{figure}

Os resultados revelam três achados críticos. \emph{Primeiro}, a taxa de
``código ausente'' (modelos que não geram nenhum bloco de código) é alta:
5/9 para mimo, 6/9 para big-pickle e 2/9 para nemotron. Isso indica que
muitos modelos respondem em linguagem natural quando solicitados a gerar
código, ou geram blocos de formato incorreto. \emph{Segundo}, entre os
códigos gerados, a taxa de erro de runtime é substancial: 2/4 para mimo,
2/3 para big-pickle e 4/7 para nemotron --- os modelos geram código que
\emph{não executa}. \emph{Terceiro}, a comparação R503 vs R506 é
iluminadora: mimo cai de 0,222 (texto) para 0,111 (código); big-pickle
cai de 0,111 para 0,000; nemotron permanece em 0,222. A capacidade de
gerar código executável correto é \emph{consistente e mensuravelmente
inferior} à de gerar texto que parece conter a resposta.

O nemotron se destaca por gerar mais código (7/9 problemas, média 317\,B)
mas com alta taxa de erro de runtime (4/7), indicando ``confiança
descalibrada'' --- o modelo gera código elaborado que falha na execução.
O big-pickle gera pouco código (3/9) mas quando gera, também falha (2/3
runtime errors). O mimo gera código intermediário (4/9) com taxa de erro
moderada (2/4).
% 03-discussao-pt.tex — Discussão
\section{Discussão}
\label{sec:discussion}

\subsection{Raio-x da arquitetura (camadas C0--C5)}
\label{sec:xray}
O experimento isola quais camadas de orquestração realmente importam.
\begin{itemize}
  \item \textbf{C0 --- Runner}: \texttt{opencode run} funciona (32/38 chamadas ok, 84\%).
  \item \textbf{C1 --- Formatos de scaffold}: a condição apenas contexto
  falhou (nemotron timeout, saídas vazias) --- formatação sozinha é
  insuficiente sem scaffold de raciocínio.
  \item \textbf{C2 --- MASWOS de 16 estágios}: decisivo; entre os modelos
  que resolveram, 100\% de sucesso no primeiro trial sob C2 vs.\ 0\% sob C0.
  \item \textbf{C3 --- Verificação determinística}: substituiu o grader com
  vazamento, expondo sensibilidade de serialização e respostas parciais.
  \item \textbf{C4 --- Validação cruzada}: a taxa de problema único de 0,98
  (mimo) colapsou para 0,22 em 9 problemas --- a maior correção do estudo.
  \item \textbf{C5 --- Trust Engine / Token Economy}: scores de confiança
  anteriores são atualizados por ciclo; a evidência R503 (big-pickle 1/9)
  reduziria seu trust score antes de qualquer delegação de roteamento.
  \item \textbf{C6 --- Governança Epistemológica (R504)}: Evidence Guard
  impede que memória/policy promovam OBSERVED; Brier/ECE medem calibração;
  drift detection bloqueia readiness; Ledger SHA-256 garante integridade
  auditável.
  \item \textbf{C7 --- Ponte Hermes (R505)}: contratos opcionais de
  interoperabilidade com runtime Hermes, sem autoridade epistemológica;
  bridge inerte sem transporte.
\end{itemize}

\subsection{Comparação com a literatura}
\label{sec:comparacao}
LLMs gratuitos atingem taxas substancialmente inferiores às reportadas em
sistemas pagos/altos parâmetros na mesma classe de dificuldade
\citep{hendrycks_math,cobbe_gsm8k,hendrycks_mmlu}. A comparação direta é
confundida por seleção de itens, temperatura, scaffold e número de trials;
é por isso que enfatizamos a \emph{correção} e não um ranking pontual, e
reportamos IC em vez de alegações. O efeito de scaffold espelha
Chain-of-Thought \citep{wei_cot}, Reflexion \citep{shinn_reflexion} e ReAct
\citep{yao_react}; a falha de saída vazia/raciocínio falso dos muse-spark
lembra as descontinuidades de capacidade alertadas por
\citet{schaeffer_mirage}. O passo de verificação determinística paralela ao
padrão ``uso de ferramentas''/prova formal \citep{yang_leandojo,alpha_geometry}.
A própria arquitetura é um sistema multiagente no espírito de AutoGen
\citep{wu_autogen} e MetaGPT \citep{hong_metagpt}, mas com gates SDD/TDD
auditáveis e loop metacognitivo (mecanismo Reflexion
\citep{shinn_reflexion}).

\subsection{Implicações para a calibração de confiança}
\label{sec:calibracao}
O ReversaFeynman Bridge (R504) revelou que o Trust Engine anterior
atualizava scores \emph{sem calibração formal}. A avaliação offline v3
mostra:
\begin{itemize}
  \item Brier score de 0,711 e ECE de 0,756 com os dados do R503 --- alarme
  claro de que as confianças originais (mimo 0,987; big-pickle 0,923) são
  \emph{mal calibradas} em relação às acurácias reais (0,22 e 0,11);
  \item readiness bloqueado corretamente (minRecords 30; temos apenas 9);
  \item drift detection: insufficient\_data (janelas maiores que a amostra);
  \item calibração do roteador: mimo 0,987 $\rightarrow$ 0,743; big-pickle
  0,923 $\rightarrow$ 0,558 (recomendação shadow, applied=False).
\end{itemize}

\subsection{Limitações}
\label{sec:limitacoes}
\begin{enumerate}
  \item $n=9$ problemas e $n=1$ trial por célula: poder estatístico baixo;
  IC amplos (ex.: mimo [0,03; 0,60]).
  \item Extração de respostas por regex; um modelo pode resolver mas
  formatar diferente.
  \item Variabilidade de rede/provedor: timeouts diferentes por modelo.
  \item Sem controle de temperatura para esses modelos; variância de
  single-run não controlada.
  \item O roteamento (R500) favorece capacidade demonstrada e não pode
  generalizar out-of-distribution.
  \item Variabilidade entre execuções (mimo variou de 0,44 para 0,22)
  exige múltiplos trials, não executados aqui por restrição de tempo.
\end{enumerate}

\subsection{Anti-overclaim}
\label{sec:anti-overclaim}
Seguindo a política do projeto \citep{corrigendum_internal}, jamais declaramos
``superhuman'', ``verificado'' ou ``Qualis A1'' sem validação externa.
Uma banca de 8 scanners (EXS) pontuou este manuscrito 47,5/100
(requires\_refinement); reportamos essa pontuação honestamente. Todos os 15
DOIs referenciados resolvem com HTTP 200 \citep{dois_validados}.
% 04-conclusao-pt.tex — Conclusão
\section{Conclusão}
\label{sec:conclusion}

Apresentamos um pipeline multiagente auditável e em evolução para avaliação
de LLMs \emph{gratuitos e reais} em matemática genuína de olimpíadas,
desenvolvido nos ciclos R497--R505 com uma contra-prova de validação cruzada
(R503) executada após reinício do host.

\emph{Respostas à pergunta de pesquisa.}
\begin{enumerate}
\item \emph{Modelos gratuitos podem ser orquestrados e scaffolddados}: o
scaffold MASWOS de 16 estágios (C2) foi decisivo --- de 0\% a 100\% de
sucesso no primeiro trial entre os resolveres --- enquanto condições brutas
e com apenas contexto falharam com saídas vazias. O orquestrador, o
roteamento (R500) e a verificação determinística (R499) formam um stack
reproduzível e auditável.
\item \emph{Honestidade de benchmarking importa}: o ranking de problema
único (mimo 0,98 $>$ big-pickle 0,92 $>$ nemotron 0,71) \emph{não}
sobrevive à validação cruzada. Taxas corrigidas são mimo 0,22 (IC 95\%
[0,03; 0,60]), big-pickle 0,11 [0,00; 0,48], nemotron 0,22 [0,03; 0,60];
McNemar pareados não-significativos; $\kappa$ mimo$\times$big-pickle 0,609
(concordância substancial --- ambos acertam o mesmo problema). Toda inferência
é reportada com intervalos de confiança, testes pareados e distribuições de
atraso, e limitações são explícitas.
\item \emph{Variabilidade entre execuções é real}: mimo variou de 0,44 para
0,22 entre duas corridas independentes --- reforçando que benchmarks de LLMs
exigem múltiplos trials por célula.
\item \emph{Artefatos metodológicos foram expostos e corrigidos}: o vazamento
do arsenal, a sensibilidade de serialização e o ``raciocínio falso'' dos
muse-spark; a hipótese do estado do host para big-pickle foi refutada.
\item \emph{Governança epistemológica funciona}: o Brier score de 0,711 e
ECE de 0,756 provam que o Trust Engine anterior estava mal calibrado; a
Avaliação Offline v3 (R504) bloqueia readiness corretamente; o Ledger
SHA-256 assegura integridade auditável.
\end{enumerate}

\emph{Trabalhos futuros.} Repetições com $n \gg 9$ problemas e múltiplos
trials por célula; ganchos de controle de temperatura; divisão treino/
teste entre desenvolvimento e avaliação; replicação independente por
terceiros; e integração da camada de verificação determinística como gate
do Trust Engine para delegação automática.

\textbf{Disponibilidade.} Todas as especificações (SPEC-935-R497..R505),
arquivos de teste, dados brutos, estatísticas e figuras estão armazenados
no repositório do OpenCode Ecosystem Core para auditoria e reprodução.

\bibliography{refs_pt}
\end{document}